\section{Supports with at most two elements}
\label{sec:small}

This section deals with the supports with at most two elements. By
\autoref{thm:palindromy} a globally periodic equation of the form
\eqref{eq:familya} has a palindromic support, and by \autoref{lem:unfold} it
suffices to treat primitive supports. A palindromic support with at most two
elements is $\emptyset$, $\{k/2\}$ or $\{j,k-j\}$ with $1\le j<k/2$, since
$k/2$ is the only index paired with itself. Primitivity then leaves three
cases:
\begin{itemize}
\item $S=\emptyset$, which is primitive only at $k=1$. There the reciprocal
form is $x_{n+1}x_n=c_0$ with $c_0>0$, which is of known type, so there is
nothing to prove.
\item $S=\{k/2\}$, which is primitive only at $k=2$: the support $S=\{1\}$
of order two (\autoref{thm:order2}).
\item $S=\{j,k-j\}$ with $\gcd(j,k)=1$. The case $k=3j$ forces $j=1$, and
is the support $\{1,2\}$ of order three (\autoref{thm:order3}); every other
case is settled by \autoref{thm:twoterm}.
\end{itemize}

At orders two and three with $c_0>0$ neither the linearisation nor the
min-plus recursion can decide the question: at order two the characteristic
polynomial $\lambda^2-\eps_1\lambda+1$ has roots of unity whenever
$\eps_1=2\cos(2\pi q/p)$ for integers $q$ and $p\ge1$, and with $\theta=1$ the min-plus recursion is itself
periodic, with period $5$ at order two and $8$ at order three
\citep[Theorem~9]{CsornyeiLaczkovich2001}. What decides it is the second half
of \autoref{thm:P}: once round a cycle of the min-plus recursion the leading
coefficients are multiplied by a fixed factor, which must be $1$.

\subsection{Order two}

At $k=2$ with $S=\{1\}$ the equation \eqref{eq:familya} reads
\begin{equation}
x_{n+2}x_n=c_0+c_1x_{n+1},\qquad c_1>0,\quad c_0\ge0.
\label{eq:small-order2}
\end{equation}
The substitution $x_n=c_1y_n$ turns it into
\begin{equation}
y_{n+2}y_n=\alpha+y_{n+1},\qquad \alpha=\frac{c_0}{c_1^2}\ge0,
\label{eq:small-lyness}
\end{equation}
and it preserves positivity and periodicity of solutions. For $\alpha=1$
this is the Lyness equation \citep{Lyness1942}, of period $5$; for
$\alpha=0$ it has period $6$. The following theorem is known:
\citet{KocicLadasRodrigues1993} showed that the periods $5$ and $6$ occur only
for $\alpha=1$ and $\alpha=0$, and the theorem itself was proved by
\citet{BarbeauGelbordTanny1995} on the whole real plane, by
\citet{CsornyeiLaczkovich2001} with complex coefficients, and by
\citet[Theorem~3.1]{CimaGasullManosas2004}. I sketch the proof, which is
essentially that of \citet[Lemma~5]{CsornyeiLaczkovich2001} and of a remark of
\citet[Section~3]{BarbeauGelbordTanny1995}, because it is the model for order
three, where the details are given. The formalisation (\autoref{app:lean})
contains a complete proof along the lines of the sketch.

\begin{theorem}[order two]
\label{thm:order2}
If \eqref{eq:small-order2} is globally periodic, then $c_0=0$ or
$c_0=c_1^2$.
\end{theorem}

\begin{proof}[Sketch of proof]
Let $c_0>0$, so that $\alpha>0$ and $\theta=1$. The min-plus recursion
$e_{n+2}+e_n=\min(0,e_{n+1})$ started from $(e_0,e_1)=(1,0)$ runs through
the $5$-cycle $1,0,-1,-1,0$. Following the leading coefficients
\eqref{eq:trop-lead} of \autoref{thm:leading} once round this cycle, with
$c_1=1$ and constant term $\alpha$, multiplies $C_{5s+1}$ by $\alpha$: the first
two steps give $C_{5s+4}=1/C_{5s+1}$, and at $n=5s+4$ the constant slot alone
attains the minimum, so $C_{5s+6}=\alpha/C_{5s+4}=\alpha\,C_{5s+1}$. Hence
$C_{5s+1}=\alpha^s$. By \autoref{thm:P} the period $p$ is a multiple of $5$,
and $C_{p+1}=C_1$ gives $\alpha^{p/5}=1$, that is $\alpha=1$.
\end{proof}

\subsection{Order three}

At $k=3$ with $S=\{1,2\}$, \autoref{thm:palindromy} gives $c_1=c_2$ for a
globally periodic equation, so the equation is
\begin{equation}
x_{n+3}x_n=c_0+c_1(x_{n+1}+x_{n+2}),\qquad c_1>0,\quad c_0\ge0.
\label{eq:small-order3}
\end{equation}
The substitution $x_n=c_1y_n$ turns it into
\begin{equation}
y_{n+3}y_n=\alpha+y_{n+1}+y_{n+2},\qquad\alpha=\frac{c_0}{c_1^2}\ge0.
\label{eq:small-todd}
\end{equation}
For $\alpha=1$ this is the period-$8$ equation of \eqref{eq:fivefamilies}.
The following theorem is also known. With complex coefficients it follows from
Theorems~1 and~4 of \citet{CsornyeiLaczkovich2001}, and with nonnegative
coefficients it is the case $m=1$ of \citet[Theorem~4.1]{CimaGasullManosas2004};
\citet{KocicLadasRodrigues1993} had shown that the period $8$ occurs exactly
when $c_0=c_1^2$. For $c_0>0$
the proof below is that of \citet[Lemma~5(ii)]{CsornyeiLaczkovich2001}, written
in the language of \autoref{thm:leading}.

\begin{theorem}[order three]
\label{thm:order3}
If $x_{n+3}x_n=c_0+c_1x_{n+1}+c_2x_{n+2}$, with $c_1,c_2>0$ and $c_0\ge0$,
is globally periodic, then $c_1=c_2$ and $c_0=c_1^2$. In particular
$c_0>0$.
\end{theorem}

\begin{proof}
Suppose the equation is globally $p$-periodic with $p\ge1$. By
\autoref{thm:palindromy}, $c_1=c_2$, so it has the form
\eqref{eq:small-order3}. The two values of $\theta$ are handled separately.

First suppose $c_0>0$, so that \eqref{eq:small-todd} is globally
$p$-periodic with $\alpha>0$. The argument follows the sketch for order two. The min-plus
recursion is $e_{n+3}+e_n=\min(0,e_{n+1},e_{n+2})$, and the solution with initial
window $(e_0,e_1,e_2)=(1,0,-1)$ is $8$-periodic, with
\[
\begin{array}{c|cccccccc}
n\bmod8 & 0 & 1 & 2 & 3 & 4 & 5 & 6 & 7\\
\hline
e_n & 1 & 0 & -1 & -2 & -2 & -1 & 0 & 1\\
A_n & \{2\} & \{2\} & \{1,2\} & \{1\} & \{1\} & \{0,1\} & \{0\} & \{0,2\}
\end{array}.
\]
The period is exactly $8$, so \autoref{thm:P} gives $p=8r$ with $r\ge1$. If
$C_{8s}=\xi$, $C_{8s+1}=\eta$ and $C_{8s+2}=\zeta$, the eight steps
$n=8s,\dots,8s+7$ of \eqref{eq:trop-lead} give in turn
\begin{gather*}
C_{8s+3}=\frac{\zeta}{\xi},\quad
C_{8s+4}=\frac{\zeta}{\xi\eta},\quad
C_{8s+5}=\frac{\eta+1}{\xi\eta},\quad
C_{8s+6}=\frac1\eta,\quad
C_{8s+7}=\frac{\eta+1}{\zeta},\\
C_{8s+8}=\frac{\xi(\alpha\eta+1)}{\eta+1},\quad
C_{8s+9}=\alpha\eta,\quad
C_{8s+10}=\alpha\zeta .
\end{gather*}
For example, at $n=8s+2$ both taps attain, so
$C_{8s+5}C_{8s+2}=C_{8s+3}+C_{8s+4}$. Since $C_1=1$, induction gives
$C_{8s+1}=\alpha^s$, and \autoref{thm:P} gives $\alpha^r=C_{p+1}=C_1=1$. Hence
$\alpha=1$, that is $c_0=c_1^2$.

Now suppose $c_0=0$, so the equation is $x_{n+3}x_n=c_1(x_{n+1}+x_{n+2})$.
Its positive equilibrium is $x^*=2c_1$, so $\eps_1=\eps_2=\tfrac12$ in
\eqref{eq:red-ej}, and the
characteristic polynomial of \autoref{thm:spectral} is
\[
\chi(\lambda)=\lambda^3-\tfrac12\lambda^2-\tfrac12\lambda+1
=(\lambda+1)\bigl(\lambda^2-\tfrac32\lambda+1\bigr).
\]
By \autoref{thm:spectral} every root of $\chi$ is a $p$-th root of unity.
Let $\lambda$ be a root of the quadratic factor. Its discriminant is
negative, so its two roots are $\lambda$ and $\bar\lambda$, and
$\lambda\bar\lambda=1$; hence $\lambda+\lambda^{-1}=\tfrac32$. But $\lambda$
is a root of unity, hence an algebraic integer, and so is $\lambda^{-1}$;
so $\tfrac32$ would be a rational algebraic integer, that is, an integer.
This contradiction shows that $c_0=0$ is impossible.
\end{proof}

\subsection{Two-term supports}

The remaining small supports are the two-term supports $\{j,k-j\}$ at
every order.

\begin{theorem}[two-term supports]
\label{thm:twoterm}
Let $1\le j$, $2j<k$, $\gcd(j,k)=1$ and $k\ne3j$, and let $S=\{j,k-j\}$.
Then no equation \eqref{eq:familya} with support $S$ is globally periodic,
for either value of $\theta$.
\end{theorem}

Note that $\gcd(S\cup\{k\})=\gcd(j,k)$, so the hypothesis $\gcd(j,k)=1$ is
exactly primitivity of $S$; and for primitive $S$ the condition $k\ne3j$
only excludes $(k,j)=(3,1)$, the order-three support of
\autoref{thm:order3}.

Throughout this subsection, for $t\in\C$ I write
\begin{equation}
\chi_t(\lambda)=\lambda^k+1-t\,(\lambda^j+\lambda^{k-j}),
\end{equation}
which is the characteristic polynomial of \autoref{thm:spectral} for the
support $\{j,k-j\}$ with $\eps_j=\eps_{k-j}=t$. Let $\eta_1,\dots,\eta_k$
be its roots, listed with multiplicity, and let
$s_r=\sum_{i=1}^k\eta_i^r$ be their power sums.

The first step is to compute a few power sums. Because $\chi_t$ has only
four terms, Newton's identities collapse to a short recursion.

\begin{lemma}[power sums]
\label{lem:small-newton}
Let $1\le j$ and $2j<k$. Then
\begin{enumerate}
\item for $1\le r<k-j$, $s_r=j\,t^{r/j}$ if $j\mid r$, and $s_r=0$
otherwise;
\item $s_{k-j}=(k-j)\,t+t\,s_{k-2j}$;
\item $s_k=t\,(s_{k-j}+s_j)-k$.
\end{enumerate}
\end{lemma}

\begin{proof}
Write $\chi_t(\lambda)=\lambda^k+a_1\lambda^{k-1}+\dots+a_k$. Since
$1\le j<k-j<k$, the only nonzero coefficients are $a_j=a_{k-j}=-t$ and
$a_k=1$. Newton's identities
$s_r+a_1s_{r-1}+\dots+a_{r-1}s_1+r\,a_r=0$, for $1\le r\le k$, therefore
read
\[
s_r=t\,s_{r-j}\,[j<r]+t\,s_{r-(k-j)}\,[k-j<r]
+r\,t\,[r=j]+r\,t\,[r=k-j]-k\,[r=k],
\]
where $[\,\cdot\,]$ is $1$ if the condition holds and $0$ otherwise. For
$r<k-j$ only the first and third terms can be present, so $s_r=t\,s_{r-j}$
for $j<r<k-j$, $s_j=jt$ and $s_r=0$ for $r<j$; part 1 follows by
induction on $r$. Parts 2 and 3 are the cases $r=k-j$ and $r=k$, using
$k-(k-j)=j$.
\end{proof}

The power sums replace a delicate analysis of where the roots of $\chi_t$
leave the unit circle: if every $|\eta_i|=1$ then $|s_r|\le k$ for every $r$,
and for a suitable $r$ the power sum $s_r$ is an explicit function of $t$
that grows too fast.

\begin{lemma}[interval bound]
\label{lem:small-bound}
Let $1\le j$, $2j<k$, $\gcd(j,k)=1$ and $k\ne3j$, and let $t\in\C$. If
every root of $\chi_t$ has modulus $1$, then $t$ is real, $|t|\le2$ and
$t\le\sqrt2$.
\end{lemma}

\begin{proof}
First, $t$ is real. The polynomial $\chi_t$ is self-reciprocal,
$\lambda^k\chi_t(1/\lambda)=\chi_t(\lambda)$, and $\chi_t(0)=1$, so
\[
\chi_t(\lambda)=\lambda^k\chi_t(1/\lambda)=\prod_i(1-\eta_i\lambda)
=\prod_i(-\eta_i)\prod_i\bigl(\lambda-\eta_i^{-1}\bigr)
=\prod_i\bigl(\lambda-\eta_i^{-1}\bigr),
\]
since $\prod_i(-\eta_i)=\chi_t(0)=1$. So the inverses $\eta_i^{-1}$ are the roots
of $\chi_t$ again, with multiplicity. As $|\eta_i|=1$, $\eta_i^{-1}=\bar\eta_i$,
and therefore $s_j=\sum_i\eta_i^{-j}=\sum_i\bar\eta_i^{\,j}=\bar s_j$. By
\autoref{lem:small-newton}, $s_j=jt$, so $t=\bar t$ is real.

Second, $|s_r|\le\sum_i|\eta_i|^r=k$ for every $r$. Suppose $j\ge2$. Then
$0<k-2j<k-j$, and $j\nmid k-2j$, since otherwise $j\mid k$ and
$j=\gcd(j,k)=1$. By \autoref{lem:small-newton},
$s_{k-2j}=0$, $s_{k-j}=(k-j)t$, $s_j=jt$ and
\[
s_k=t\bigl((k-j)t+jt\bigr)-k=k\,(t^2-1).
\]
So $k|t^2-1|\le k$, whence $t^2\le2$, and $|t|\le\sqrt2$.

Suppose instead $j=1$. Then $k>2$ and $k\ne3$, so $k\ge4$, and
\autoref{lem:small-newton} gives $s_{k-2}=t^{k-2}$ and
$s_{k-1}=(k-1)t+t\,s_{k-2}=(k-1)t+t^{k-1}$. From the first,
$|t|^{k-2}\le k\le2^{k-2}$ (as $k\ge4$), so $|t|\le2$. From the second, if
$t>1$ then $s_{k-1}>(k-1)+1=k$, which is impossible; so $t\le1\le\sqrt2$.
\end{proof}

The last ingredient is Kronecker's theorem \citep{Kronecker1857}: if $L$ is
a number field and $\zeta\in L$ is a nonzero algebraic integer with
$|\varphi(\zeta)|\le1$ for every field embedding $\varphi:L\to\C$, then
$\zeta$ is a root of unity; for a modern account see
\citet[Theorem~1.5.9]{BombieriGubler2006}.

\begin{proof}[Proof of \autoref{thm:twoterm}]
Suppose that an equation \eqref{eq:familya} with support $S=\{j,k-j\}$ is
globally $p$-periodic. By \autoref{thm:palindromy}, $c_j=c_{k-j}$, so
$\eps_j=\eps_{k-j}=:\eps$, and \eqref{eq:red-ej} gives
$2\eps=\eps_j+\eps_{k-j}\le1$. Thus
\[
0<\eps\le\tfrac12,
\]
and by \autoref{thm:spectral} every root $\eta_i$ of $\chi_\eps$ is a
$p$-th root of unity. The idea of the proof is that $\eps=\zeta+\zeta^{-1}$
for a root of unity $\zeta$. The bound $0<\eps\le\tfrac12$ forces $\zeta$ to
have order at least $9$, and then some conjugate of $\eps$ is larger than
$\sqrt2$, which \autoref{lem:small-bound} forbids.

\emph{Step 1: $\eps$ is $\zeta+\zeta^{-1}$ for a root of unity $\zeta$.}
The coefficient of $\lambda^{k-j}$ in
$\chi_\eps=\prod_i(\lambda-\eta_i)$ is $-\eps$, and it is also $(-1)^j$
times the $j$-th elementary symmetric function of $\eta_1,\dots,\eta_k$.
Roots of unity are algebraic integers, so $\eps$ is an algebraic integer. Let $\zeta$ be a
root of $z^2-\eps z+1$; then $\zeta$ is an algebraic integer (a root of a
monic polynomial whose coefficients are algebraic integers), $\zeta\ne0$,
and $\eps=\zeta+\zeta^{-1}$. Let $L=\Q(\zeta,\eta_1,\dots,\eta_k)$,
a number field containing $\eps$. Let $\varphi:L\to\C$ be an embedding.
Applying $\varphi$ to the coefficients of
$\chi_\eps=\prod_i(\lambda-\eta_i)$ gives
$\chi_{\varphi(\eps)}=\prod_i(\lambda-\varphi(\eta_i))$, and each
$\varphi(\eta_i)$ is again a $p$-th root of unity, so of modulus $1$. By
\autoref{lem:small-bound}, $r=\varphi(\eps)$ is real with $|r|\le2$. Now
$\varphi(\zeta)$ is a root of $z^2-rz+1$, a real quadratic with
discriminant $r^2-4\le0$; its two roots are therefore complex conjugates
(or equal to $\pm1$) with product $1$, so $|\varphi(\zeta)|=1$. By
Kronecker's theorem, $\zeta$ is a root of unity.

\emph{Step 2: the order of $\zeta$.} Let $n_0$ be the order of $\zeta$.
Since $|\zeta|=1$, $\zeta^{-1}=\bar\zeta$, so
$\eps=2\cos(2\pi b/n_0)$ for some integer $b$ coprime to $n_0$, and I may
take $0\le b\le n_0/2$. Since $\cos$ is decreasing on $[0,\pi]$, with
$\cos(5\pi/12)=(\sqrt6-\sqrt2)/4>\tfrac14$ and $\cos(\pi/2)=0$, the
condition $0<\eps\le\tfrac12$ forces
\[
\frac5{24}<\frac b{n_0}<\frac14,
\]
that is, $5n_0<24b<6n_0$. If $n_0\le8$ then $24b<48$, so $b=1$ (as $b=0$ is
excluded), and $5n_0<24<6n_0$ has no integer solution. Hence $n_0\ge9$.

\emph{Step 3: a large conjugate.} The cyclotomic polynomial
$\Phi_{n_0}$ is irreducible over $\Q$, so it is the minimal polynomial of
$\zeta$, and $\e^{2\pi i/n_0}$ is another of its roots. There is therefore a
field embedding of $\Q(\zeta)$ into $\C$ sending $\zeta$ to $\e^{2\pi i/n_0}$,
and since $\zeta\in L$ it extends to an embedding $\varphi:L\to\C$. Then
\[
\varphi(\eps)=\varphi(\zeta)+\varphi(\zeta)^{-1}
=\e^{2\pi i/n_0}+\e^{-2\pi i/n_0}=2\cos(2\pi/n_0),
\]
and, as in Step 1, \autoref{lem:small-bound} gives
$2\cos(2\pi/n_0)\le\sqrt2$. But $n_0\ge9$ gives $0<2\pi/n_0<\pi/4$, so
$2\cos(2\pi/n_0)>2\cos(\pi/4)=\sqrt2$. This contradiction proves the
theorem.
\end{proof}
